\section{Introduction}
\label{sec:introduction}

Real-time multi-class visual perception is a foundational prerequisite for safe, efficient operations in modern automated warehouse logistics. Contemporary smart warehousing environments increasingly integrate heterogeneous fleets of Autonomous Mobile Robots (AMRs), Automated Guided Vehicles (AGVs), forklifts, and human workers operating in shared, high-density workspaces. Accurate, low-latency detection of moving machinery, personnel, pallets, and goods is essential to prevent catastrophic collisions, optimize routing throughput, and enable dynamic inventory tracking.

Deploying deep neural vision detectors in industrial logistics facilities imposes severe operational constraints:
\begin{enumerate}
    \item \textbf{Severe Perspective and Scale Variation:} Fixed overhead surveillance cameras and mobile robot-mounted sensors observe targets ranging from expansive industrial forklifts ($>90{,}000\text{ px}^2$) to distant personnel and inventory markers ($<800\text{ px}^2$) along long aisles.
    \item \textbf{Dense Stacking and Pervasive Occlusions:} Repetitive pallet racks and high-density storage geometries cause severe inter-object occlusions, requiring long-range contextual visual cues to resolve ambiguity.
    \item \textbf{Embedded Edge Hardware Limits:} Mobile logistics robots and distributed intelligent edge cameras are constrained by tight thermal envelopes, limited battery budgets, and restricted onboard memory (e.g., NVIDIA Jetson Orin Nano / Xavier platforms), precluding the deployment of computationally heavy vision transformers.
\end{enumerate}

\subsection{Limitations of Existing Detectors}
Convolutional Neural Networks (CNNs) have established the dominant foundation for embedded visual perception, epitomized by the Ultralytics YOLO family~\cite{ultralytics2024yolo11}. Due to localized receptive fields and strong translational equivariance, lightweight CNN detectors such as YOLO11s deliver outstanding throughput and hardware cache locality. However, pure convolutional architectures inherently struggle to aggregate global spatial context across extensive warehouse aisles without stacking dozens of layers, which increases parameter latency.

Conversely, Vision Transformers (ViTs) and self-attention mechanisms offer unrestricted global receptive fields. Nevertheless, the standard multi-head self-attention mechanism exhibits quadratic computational and memory complexity $\mathcal{O}(N^2)$ with respect to spatial token length $N = H \times W$. Even pointwise spatial attention mechanisms (such as \texttt{C2PSA} in YOLO11) introduce substantial parameter overhead and memory access costs (MAC) that challenge embedded edge processors.

\subsection{Selective State Space Models (Mamba)}
Recently, Selective State Space Models (SSMs), notably Mamba~\cite{gu2023mamba}, Visual Mamba (VMamba)~\cite{liu2024vmamba}, and Vision Mamba (Vim)~\cite{zhu2024visionmamba}, have emerged as a compelling alternative for visual sequence modeling. By parameterizing continuous state transitions with dynamic, input-conditioned selective gating, SSMs achieve global receptive field coverage with strictly linear complexity $\mathcal{O}(N)$ in memory and computation. Despite promising results in general image classification and medical segmentation, the integration of 2D selective state-space mechanisms into lightweight, multi-scale anchor-free object detectors under strict edge constraints remains an underexplored frontier.

\subsection{Research Gap and Questions}
Existing visual Mamba research has largely focused on heavy, isotropic backbones or server-grade compute platforms, leaving a distinct gap: \textit{How can 2D selective state-space modeling be surgically integrated into lightweight real-time object detectors to capture global warehouse context without degrading edge deployability or inflating parameter counts?}

To address this challenge, we investigate three fundamental research questions:
\begin{itemize}
    \item \textbf{RQ1 (Accuracy--Efficiency Trade-off):} Can replacing the deepest self-attention bottleneck in YOLO11s with a 2D selective state-space block maintain or improve multi-class detection accuracy while reducing model parameters and FLOPs?
    \item \textbf{RQ2 (Optimization and Transfer Dynamics):} How does the insertion of an uncalibrated state-space block affect feature propagation across a pretrained convolutional backbone, and what optimization protocol stabilizes convergence?
    \item \textbf{RQ3 (Edge Viability and Latency Realism):} What is the practical latency profile of 2D selective scans on resource-constrained host processors, and how do sequential recurrences compare to hardware-accelerated parallel scans?
\end{itemize}

\subsection{Research Contributions}
\label{sec:contributions}
This work delivers three primary contributions:
\begin{itemize}
    \item \textbf{Hybrid Warehouse Architecture (YOLO11s-SS2D):} We introduce a hybrid architecture that replaces the parameter-heavy self-attention bottleneck (\texttt{C2PSA}) in YOLO11s with a 4-direction 2D Selective State Space block (\texttt{WarehouseSS2DBlock}) at Stage $P_5$ (and optionally $P_4$). This reduces model parameters by \textbf{$6.72\%$} ($8.80\text{ M}$ vs. $9.43\text{ M}$) and cuts Stage 10 attention parameters by \textbf{$63.9\%$} ($357{,}632$ vs. $990{,}976$) at $21.2\text{ GFLOPs}$, using zero-initialized residual projections to preserve identity mapping at initialization.
    \item \textbf{Zero-Leakage Multi-Class Warehouse Benchmark:} Sourced from \textbf{NVIDIA PhysicalAI-SmartSpaces}, we construct a cross-facility evaluation split ($300$ training frames with $12{,}959$ instances across Warehouse\_000 vs. $100$ validation frames with $3{,}696$ instances across Warehouse\_015) guaranteeing disjoint backgrounds and camera views to eliminate spatial-temporal leakage.
    \item \textbf{Optimization Dynamics and Safety Gains:} We diagnose and resolve the detection head cold-start failure by preserving class-agnostic bounding-box anchors. Across controlled 20-epoch training, the hybrid model delivers a \textbf{$+6.2\%$ boost in forklift detection} ($76.3\%\text{ vs. } 70.1\%\text{ mAP}_{50}$) and increases macro recall by $+7.30\%$, while uncovering the high-variance sample boundary condition of rare tail classes ($13$ \texttt{pallet\_truck} validation instances).
    \item \textbf{Physical Silicon Verification and Execution Profile:} We compile the pure PyTorch model to native ONNX without custom C++/CUDA plugins, deploying it directly onto a physical \textbf{Raspberry Pi 4 Model B} (Quad-Core ARM Cortex-A72 @ 1.5\,GHz). Steady-state execution consumes \textbf{$371.9\text{ MB}$ RSS} ($62.8\%$ memory headroom under $1.0\text{ GB}$). We analyze why single-threaded CPU sequential recurrence achieves approximately $0.72\text{ FPS}$ ($1398\text{ ms}$), providing the engineering rationale for parallel prefix-scan acceleration on edge GPUs and NPUs.
\end{itemize}

\subsection{Paper Organization}
\label{sec:organization}
The remainder of this paper is structured as follows: Section~\ref{sec:related_work} surveys lightweight object detection and visual state-space literature. Section~\ref{sec:methodology} introduces the mathematical formulation of the \texttt{WarehouseSS2DBlock} and hybrid architecture. Section~\ref{sec:experimental_setup} details the benchmark dataset, zero-leakage split, and experimental protocol. Section~\ref{sec:results} presents comparative empirical findings, ablation analyses, and latency profiles. Section~\ref{sec:conclusion} concludes with implications for edge robotics deployment.
